% =====================================================================
% CHAPTER 9: RESULT & ACCURACY
% =====================================================================
\chapter{Result \& Accuracy}
\label{ch:results}

\section{Evaluation Methodology}
\label{sec:eval-method}

TrackVision's tracking performance was evaluated through the following evidence-based methods:

\begin{enumerate}
    \item \textbf{Implementation-Level Unit Tests:} 49 Vitest tests validating Kalman filtering, Hungarian assignment, ByteTrack lifecycle, appearance fusion, class-aware NMS, and store history pruning.
    \item \textbf{Runtime Microbenchmarks:} JavaScript-side pipeline timing (NMS + ByteTrack association + Kalman predict/update) measured over 1000 frames in Node.js 22.
    \item \textbf{Telemetry Collection:} Per-frame latency breakdown (preprocess, inference, postprocess, tracking, render) and FPS recorded during manual qualitative testing.
    \item \textbf{Qualitative Observation:} Manual verification of identity continuity, occlusion handling, track merging, and responsiveness using live camera and sample video.
    \item \textbf{Session Initialization Timing:} Model loading and ONNX session creation times measured on Windows 11 / Chrome with WebGPU.
    \item \textbf{Bundle Size Analysis:} Production build artifact sizes measured from \texttt{dist/}.
\end{enumerate}

\vspace{0.3cm}

\noindent \textbf{Critical Limitation:} \textbf{The current implementation does not compute standardized MOT benchmark metrics such as MOTA, IDF1, HOTA, ID switches, false positives, or false negatives.} No formal benchmark dataset (MOT17, MOT20, KITTI, TAO, DanceTrack) is used, and no ground-truth annotations are available. The evaluation therefore focuses on engineering validation and computational performance rather than benchmark-oriented MOT accuracy.

\section{Implementation Validation: Unit Tests}
\label{sec:unit-tests}

All 49 unit tests pass (\texttt{npm run test}):

\begin{table}[htbp]
\centering
\caption{Unit test coverage by module}
\label{tab:tests}
\begin{tabular}{@{}lr@{}}
\toprule
\textbf{Module} & \textbf{Tests} \\
\midrule
Kalman Filter (\texttt{kalman.test.ts}) & 9 \\
Hungarian/Greedy Matching (\texttt{matching.test.ts}) & 8 \\
ByteTrack Lifecycle + Appearance (\texttt{tracker.test.ts}) & 20 \\
Worker Utils / NMS (\texttt{workerUtils.test.ts}) & 7 \\
Zustand Store History (\texttt{store.test.ts}) & 5 \\
\midrule
\textbf{Total} & \textbf{49} \\
\bottomrule
\end{tabular}
\end{table}

\begin{itemize}
    \item \textbf{Kalman:} Predict/update correctness, covariance propagation, matrix inversion, velocity clamping, state reset.
    \item \textbf{Matching:} Hungarian optimal assignment, gating, rectangular padding, greedy fallback, edge cases (empty inputs).
    \item \textbf{ByteTrack:} TENTATIVE$\to$CONFIRMED$\to$LOST$\to$REMOVED transitions, confirmation thresholds (hits$\ge$3 or confirmedHits$\ge$2), occlusion recovery (re-association after 3 missed frames), ID persistence during smooth motion, separate IDs for spatially distinct objects, low-score detection handling, label stabilization via majority vote.
    \item \textbf{Appearance Fusion:} EMA blending ($\alpha$=0.3) with L2 re-normalization, crossing disambiguation when IoU ambiguous (appearance-dominant tuning).
    \item \textbf{Store:} 5-minute rolling window frame pruning, track metadata persistence, event logging.
\end{itemize}

\section{Pipeline Microbenchmarks}
\label{sec:microbench}

JavaScript-side tracking pipeline (NMS + ByteTrack association + Kalman predict/update) measured over 1000 frames per configuration in Node.js 22 (Vitest benchmark):

\begin{table}[htbp]
\centering
\caption{Tracking pipeline latency (JS only, no model inference)}
\label{tab:pipeline-perf}
\begin{tabular}{@{}cc@{}}
\toprule
\textbf{Detections/Frame} & \textbf{ms/frame} \\
\midrule
5  & 0.31 \\
10 & 0.51 \\
20 & 1.04 \\
\bottomrule
\end{tabular}
\end{table}

\begin{figure}[htbp]
\centering
\begin{tikzpicture}
\begin{axis}[
    width=0.7\textwidth,
    height=0.4\textwidth,
    xlabel={Detections per frame},
    ylabel={Processing time (ms)},
    title={Tracking Pipeline Latency vs. Detection Count},
    xmin=0, xmax=22,
    ymin=0, ymax=1.2,
    xtick={5,10,20},
    ytick={0,0.31,0.51,1.04},
    mark=*, mark size=3,
    grid=major,
    legend pos=north west
]
\addplot[color=trackvisionblue, thick] coordinates {
    (5,0.31) (10,0.51) (20,1.04)
};
\addlegendentry{Measured}
\end{axis}
\end{tikzpicture}
\caption{Tracking pipeline latency scales sub-linearly with detection count. Model inference (not shown) dominates end-to-end frame time.}
\label{fig:perf-chart}
\end{figure}

\vspace{0.3cm}

\noindent \textbf{Interpretation:} The JS tracking overhead is negligible ($<$ 1.1 ms/frame even at 20 detections). End-to-end frame time is dominated by model inference (YOLO/OSNet/CLIP), which is hardware-dependent (WebGPU vs. WASM).

\section{Session Initialization Times}
\label{sec:init-times}

Measured on Windows 11, Chrome (WebGPU execution provider), production build:

\begin{table}[htbp]
\centering
\caption{ONNX session initialization time (model load + compile)}
\label{tab:init-times}
\begin{tabular}{@{}lcc@{}}
\toprule
\textbf{Model} & \textbf{Size} & \textbf{Session Init} \\
\midrule
YOLOv8n (detection) & 12.8 MB & $\sim$1.8 s \\
OSNet x1.0 (ReID)   & 8.8 MB  & $\sim$1.9 s \\
YOLO-World (open)   & 418 MB  & $\sim$3.7 s \\
\bottomrule
\end{tabular}
\end{table}

\vspace{0.3cm}

\noindent \textbf{Notes:}
\begin{itemize}
    \item Initialization includes model download (if not cached), ONNX parsing, graph optimization, and execution provider compilation.
    \item WebGPU provider adds shader compilation time; WASM is faster to initialize but slower at inference.
    \item YOLO-World's large size (418 MB) dominates its initialization time.
    \item CLIP text encoder ($\sim$254 MB) loads lazily on first concept set in open mode.
\end{itemize}

\section{Bundle Sizes (Production Build)}
\label{sec:bundle}

Measured from \texttt{dist/assets/} after \texttt{npm run build} (raw, not gzipped):

\begin{table}[htbp]
\centering
\caption{Production bundle sizes}
\label{tab:bundle}
\begin{tabular}{@{}lr@{}}
\toprule
\textbf{Chunk} & \textbf{Size} \\
\midrule
Main bundle (\texttt{index.js}) & 345 KB \\
Charts chunk (\texttt{charts.js}, lazy) & 373 KB \\
Vendor chunk & 36 KB \\
YOLO Detection Worker & 404 KB \\
YOLO-World Worker & 404 KB \\
ReID Worker & 400 KB \\
Tracker Worker & 10 KB \\
CSS & 81 KB \\
\hline
ONNX Runtime WASM (\texttt{ort-wasm-simd-threaded.jsep.wasm}) & 26.2 MB \\
\bottomrule
\end{tabular}
\end{table}

\vspace{0.3cm}

\noindent The WASM binary is cached by the browser and loaded on-demand by workers. Total initial download (without optional models): \textbf{$\sim$1.7 MB} (JS + CSS + WASM).

\section{Qualitative Tracking Behavior}
\label{sec:qualitative}

Manual testing with live camera and sample video (\texttt{one-by-one-person-detection.mp4}) observed:

\begin{itemize}
    \item \textbf{Identity Continuity:} Tracks maintain stable IDs during smooth motion; re-association succeeds after brief occlusion ($\le$ 5 frames).
    \item \textbf{Occlusion Handling:} Kalman prediction interpolates up to 5 frames during occlusion; track status transitions TRACKED $\to$ LOST after 6 missed frames.
    \item \textbf{Track Merging:} When two tracks of same class overlap (IoU > 0.5) with high appearance similarity ($>$ 0.7), older track absorbs younger.
    \item \textbf{Label Stabilization:} Majority vote over 10-frame history corrects transient misclassifications.
    \item \textbf{Open-Vocabulary Detection:} YOLO-World + CLIP successfully detects user-specified concepts (e.g., ``water bottle'', ``headphones'') when full models are downloaded.
    \item \textbf{ReID Fusion:} With OSNet loaded, appearance embeddings improve crossing disambiguation; without OSNet, tracker gracefully falls back to IoU + center + class.
\end{itemize}

\section{Standardized MOT Evaluation -- Future Work}
\label{sec:future-eval}

The following standardized evaluation is \textbf{not currently implemented} but is recommended for future work:

\begin{table}[htbp]
\centering
\caption{Recommended MOT metrics for future benchmark evaluation}
\label{tab:future-metrics}
\begin{tabular}{@{}ll@{}}
\toprule
\textbf{Metric} & \textbf{Description} \\
\midrule
MOTA & Multiple Object Tracking Accuracy (detection + association) \\
IDF1 & ID F1-score (identity consistency) \\
HOTA & Higher Order Tracking Accuracy (unified detection+association) \\
ID Switches & Count of identity changes \\
FP / FN & False positives / false negatives \\
MOTP & Multiple Object Tracking Precision (localization) \\
\bottomrule
\end{tabular}
\end{table}

\vspace{0.3cm}

\noindent \textbf{Required for standardized evaluation:}
\begin{enumerate}
    \item Annotated benchmark dataset (MOT17 train/val, MOT20, or DanceTrack).
    \item Ground-truth trajectory parser (MOT format: \texttt{<frame>, <id>, <bb\_left>, <bb\_top>, <bb\_width>, <bb\_height>, <conf>, <class>, <vis>}).
    \item Evaluation pipeline computing MOTA, IDF1, HOTA per \cite{bernardin2008evaluating, luiten2021hota}.
    \item Deterministic inference seed for reproducibility.
    \item CI integration for regression tracking.
\end{enumerate}

\section{Summary of Achieved Results}
\label{sec:results-summary}

\begin{table}[htbp]
\centering
\caption{Summary of measured results (achieved)}
\label{tab:results-summary}
\begin{tabular}{@{}ll@{}}
\toprule
\textbf{Metric} & \textbf{Value} \\
\midrule
Unit tests passing & 49 / 49 \\
Pipeline latency (5 dets) & 0.31 ms/frame \\
Pipeline latency (10 dets) & 0.51 ms/frame \\
Pipeline latency (20 dets) & 1.04 ms/frame \\
YOLOv8n init (WebGPU) & $\sim$1.8 s \\
OSNet init (WebGPU) & $\sim$1.9 s \\
YOLO-World init (WebGPU) & $\sim$3.7 s \\
Main bundle size & 345 KB \\
Total JS/CSS (excl. WASM) & $\sim$1.7 MB \\
ONNX WASM binary & 26.2 MB \\
TypeScript strict check & Pass (\texttt{tsc --noEmit}) \\
\bottomrule
\end{tabular}
\end{table}

\vspace{0.3cm}

\noindent \textbf{Explicitly NOT achieved (by design):}
\begin{itemize}
    \item MOTA / IDF1 / HOTA scores
    \item ID switch count
    \item False positive / negative rates
    \item Detection mAP on COCO
    \item End-to-end FPS claims (device-dependent; model inference dominates)
\end{itemize}

\section{Conclusion on Accuracy}
\label{sec:accuracy-conclusion}

TrackVision demonstrates \textbf{correct algorithmic implementation} (validated by 49 unit tests) and \textbf{low computational overhead} (sub-millisecond JS tracking pipeline). The system achieves real-time performance on consumer hardware with WebGPU acceleration. However, \textbf{no quantitative MOT accuracy claim can be made} without benchmark dataset evaluation. The qualitative behavior confirms the tracker functions as designed: IDs persist, occlusions are handled via Kalman prediction, and appearance fusion (when available) improves crossing disambiguation.